\section{Introduction}\label{sec:introduction}
Industrial anomaly detection (AD) is typically deployed under an asymmetric data regime: normal examples are available, while abnormal examples are rare, diverse, and costly to annotate. This makes few-shot AD practically important. In this setting, an inspection system receives only $k$ normal support images for a target category, often with $k\in\{1,2,4,8\}$, and must detect both image-level and pixel-level defects at test time.

Recent methods have made this regime much more effective by using frozen visual foundation models. DINOv2~\cite{dinov22023} patch features can be stored in a normal memory bank and queried by nearest-neighbor distances, as in PatchCore-style methods and AnomalyDINO~\cite{patchcore2021,anomalydino2024}. Alternatively, frozen patch features can be compressed into a normal PCA/subspace model, as in SubspaceAD~\cite{subspacead2026}. These approaches provide strong anomaly ranking without training a large task-specific backbone.

However, ranking is not the same as reliable decision support. A raw nearest-neighbor distance or PCA residual can produce high AUROC while still being poorly calibrated: a score threshold may work on one class, but the same numerical score may not carry the same false-alarm meaning under another class, dataset, or corruption. This matters in industrial deployment because operators need to know not only which samples are suspicious, but also whether an alarm probability is trustworthy. Prior work has already shown that DINOv2-based few-shot AD can be poorly calibrated and adversarially fragile~\cite{khan2025calibration}. Reliability is therefore not a cosmetic post-processing issue; it is part of the deployment problem.

We address this gap by separating anomaly ranking from reliability estimation. The resulting questions are about the learning system itself, not about any one factory: what false-alarm behavior can a decision rule calibrated on $k$ examples of a frozen representation guarantee, at what resolution, and how does that guarantee fail when the representation's inputs shift? The ranking path remains a frozen DINOv2 PCA/subspace detector: patch residuals are aggregated into a raw anomaly score used for AUROC, AP, and localization. A separate reliability path maps this evidence into calibrated or probability-like outputs. We call the resulting framework \emph{Conformal Reliability Routing} (CRR). The main CRR route uses leave-one-image-out (LOIO) conformal calibration over the few normal support images. Each support image is scored as if held out, producing normal nonconformity values; a test image then receives a conformal p-value and a reliability score $1-p$. This view is not a supervised anomaly posterior, but it is statistically interpretable and derived without target anomaly labels.

Each ingredient we build on---DINOv2 features, PCA residuals, Platt-family calibration, and conformal anomaly detection---has an established literature, and our contribution begins after them. What that literature leaves open is a structural question about the few-shot regime itself: a conformal alarm calibrated on $k$ support images cannot fire below $\alpha=1/(k{+}1)$ at all, corruption shift can make its observed false-alarm rate anti-conservative, and a statistically valid source-category certificate can itself be vacuous when only a few categories are available. The central contribution combines an exact resolution-floor diagnosis, an empirical audit under shift, and a strict nested evaluation of a source-assisted cross-category scheme (SC3R). The strict experiment does not validate the hoped-for new-category claim: all category-certified thresholds fall back to zero. That negative result is predictable from a feasibility bound and is reported alongside, rather than replaced by, fixed-archive image sensitivity and historical empirical results.

The lesson for learning systems is general even though the testbed is industrial: finite calibration limits both the resolution of target-only alarms and the feasibility of source-category certification. The paper makes four contributions, the first two central:
\begin{itemize}
    \item We make the few-shot resolution limit of conformal alarms explicit through an attainable-alpha analysis: with $k$ supports no alarm can fire below $\alpha=1/(k{+}1)$, and empirical false-alarm rates must be read against the attainable grid, not the nominal level. Cluster-aware Monte Carlo diagnostics of normal p-values then localize when and in which direction the exchangeable-null pattern fails under corruption shift---conservative on VisA, anti-conservative on corrupted MVTec.
    \item We formulate strict nested SC3R with disjoint reference/proposal/certification categories and derive a necessary sample-size condition for its Hoeffding gate. The frozen experiment spans four $k$, five seeds, five conditions, MVTec, VisA, and MVTec-to-VisA/MPDD transfer. Category certification returns the zero threshold in all 960 gate cells because only three or four certification categories are available; image-level fixed-archive sensitivity is reported separately and never relabeled as a new-category result.
    \item We formulate a decoupled pipeline: frozen DINOv2 PCA residuals provide ranking, while a label-free layer provides Vector/Shift-Aware calibration and LOIO p-values~\cite{hennhofer2024loio} from the normal supports. Across the historical MVTec and VisA corruption studies, LOIO lowers ECE against most tested calibrators without changing the raw ranking; operational metrics, rather than ECE, carry the central interpretation.
    \item We report accuracy-storage, transfer calibration, pixel-level, robustness, prevalence-stress, and routing diagnostics, including negative findings that prevent overclaiming: CRR is not MVTec AUROC SOTA, is not adversarially robust, ECE on conformal-derived scores is strongly prevalence-sensitive, and the historical SC3R gate is not an independent certificate.
\end{itemize}

The remainder of this paper is organized as follows. Section~\ref{sec:related} positions the work against few-shot detectors, calibration, and conformal anomaly detection. Section~\ref{sec:method} presents the decoupled pipeline, the LOIO conformal route, the attainable-alpha analysis, and SC3R. Section~\ref{sec:experiments} describes datasets, protocols, and claim rules. Section~\ref{sec:results} reports calibration comparisons, uniformity tests, operational false-alarm behavior, SC3R evaluations with their pre-specified gate, and diagnostics. Section~\ref{sec:limitations} discusses limitations, and Section~\ref{sec:conclusion} concludes.
